\origpage{41}{GRAPH MINORS}
Our proof of (1.1) is basically by induction on the ``complexity'' of the
forest $H$. This is explained in detail in Section 3, and the induction argument
is performed in Sections 4 and 5. Section 2 contains some crucial lemmas
about ``grids.''

Our terminology is mostly standard, but a few terms need explanation. If
$X, Y$ are sets, $X-Y$ denotes
\begin{equation*}
\{x\colon x\in X, x\notin Y\}.
\end{equation*}
If $G$ is a graph, $V(G)$ and $E(G)$ denote its vertex-set and edge-set, respectively.
A graph is \textit{simple} if it has no loops or multiple edges. If $X\subseteq V(G)$,
$G\backslash X$ denotes the result of deleting from $G$ all the vertices in $X$, and $G\,|\,X$
denotes $G\backslash(V(G)-X)$. A graph $G$ is a \textit{subdivision} of a graph $H$ if $H$ can be
obtained from $G$ by (repeatedly) choosing a vertex of valency 2 and
contracting an edge incident with it which is not a loop. Two subgraphs of a
graph are \textit{disjoint} if they have no common vertices, and a set of subgraphs is
disjoint if every pair of its members are disjoint. Every \textit{path} has at least one
vertex, and no ``repeated'' vertices. It has an \textit{initial} and a \textit{terminal} vertex
(which are equal for a one-vertex path) called its \textit{ends}, and any other vertices
are called \textit{internal} vertices. A path \textit{avoids} $X\subseteq V(G)$ if it has no vertex in $X$.
$X\subseteq V(G)$ \textit{separates} $Y, Z\subseteq V(G)$ if no path in $G$ from $Y$ to $Z$ avoids $X$.
(Thus $Y\cap Z\subseteq X$.) We say that $X$ separates $y, z\in V(G)$ if it separates
$\{y\}, \{z\}$.

We shall require the following lemma:
\stmt
(1.7). \textit{Let $\rho$ be a vertex of a graph $G$, and let $Y_1, Y_2,..., Y_T$ be disjoint
subsets of $V(G)$, where $T\ge 1$. Suppose that for $1\delm{<}\fix{\le} i, i'\le T$,\fixnote{fix:1.7}{As printed, nothing is assumed about $Y_1$, and the conclusion fails. For the path $\rho\,a\,b$ with $T=2$, $Y_1=\{b\}$ and $Y_2=\{a\}$, the hypothesis holds, but $W_0=\{\rho,a\}$ meets $W_2=\{a,b\}$. The proof uses the case $i=1$ (in the uniqueness step, $Y_{i+1}$ with $i=0$ is $Y_1$), and the use in Section~4, through (4.5), has $1\le i,i'\le T$.} and for all
$y\in Y_{i'}$, $Y_i$ separates $y$ and $\rho$ if and only if $i\le i'$. For $1\le i<T$, let $W_i$ be
the set of all vertices of $G$ separated from $\rho$ by $Y_i$, and not separated from $\rho$
by $Y_{i+1}$. Let $W_0$ be the set of vertices not separated from $\rho$ by $Y_1$, and let
$W_T$ be the set of all vertices separated from $\rho$ by $Y_T$. Then the sets
$W_0, W_1,..., W_T$ are disjoint and partition $V(G)$, and $Y_i\subseteq W_i$ $(1\le i\le T)$, and
the sequence}
\begin{equation*}
W_0\cup Y_1, W_1\cup Y_2,..., W_{T-1}\cup Y_T, W_T
\end{equation*}
\textit{is a path-decomposition of $G$.}
\stmt
\textit{Proof}. For every vertex $v\in V(G)$, there is a unique value of $i$ with
$0\le i\le T$ such that

\hspace*{12pt}(i)\enspace either $i=0$ or $Y_i$ separates $v$ and $\rho$, and

\hspace*{12pt}(ii)\enspace either $i=T$ or $Y_{i+1}$ does not separate $v$ and $\rho$.

\noindent For it is clear that there is at least one such value of $i$; and if $i, i'$ are two
